\documentclass[11pt]{article}
\usepackage[T1]{fontenc}
\usepackage{lmodern,amsmath,amssymb,amsthm,microtype}
\usepackage[margin=1in]{geometry}
\usepackage[hidelinks]{hyperref}
\newtheorem{theorem}{Theorem}
\newtheorem{corollary}[theorem]{Corollary}
\newcommand{\R}{\mathbb R}
\newcommand{\Sg}{\mathcal S}
\title{Full stress stratifications determine a labelled graph}
\author{Alec Kriebel\thanks{Independent researcher. ORCID:
\href{https://orcid.org/0009-0001-9320-500X}{0009-0001-9320-500X}.}}
\date{12 October 2026\quad (research note)}
\begin{document}
\maketitle
\begin{abstract}
For a simple labelled graph, consider the full configuration space, including
vertex collisions, partitioned into connected components of coordinate-sign
equivalence classes of self-stress spaces. We prove that two graphs induce
identical partitions precisely when their edge sets agree. More explicitly, an
edge is recovered by asking whether its collision diagonal is a union of strata.
This answers the literal equality formulation of Karpenkov's Problem~4.
The proof combines the known singleton-support collision criterion with an
explicit positive diagonal transport along a collinear path. The comparison
convention is essential: nonisomorphic trees can have ambient-linearly equivalent
stratifications.
\end{abstract}

\section{The full partition and its reconstruction}
Fix integers $n,d\geq1$. A graph below is simple and has the common labelled
vertex set $[n]$; unused vertices are isolated. Its configuration space is
$X=(\R^d)^n$. For $P=(p_1,\ldots,p_n)$ and a graph $G$, put
\[
 W_G(P)=\left\{w\in\R^{E(G)}:
 \sum_{j:\{i,j\}\in E(G)}w_{ij}(p_j-p_i)=0\text{ for every }i\right\}.
\]
There is one real coordinate per unordered edge, so $w_{ij}=w_{ji}$.
Two configurations for the \emph{same} graph are equivalent when a homeomorphism
between their stress spaces preserves the sign of every edge coordinate,
including zero. Let $\Sg_d(G)$ be the partition of $X$ into connected components
of these equivalence classes. We retain the entire Cartesian space: collisions,
self-intersections and deficient affine spans are allowed. These are the
full-space conventions of \cite{DKS,Karpenkov}.

Equality of partitions means equality as collections of subsets of this
\emph{fixed labelled space}. We do not directly compare stress spaces belonging
to different graphs. Write
$\Delta_{ij}=\{P\in X:p_i=p_j\}$ for a collision diagonal.

\begin{theorem}\label{thm:main}
For every simple labelled graph $G$ and every $i\ne j$,
\[
 \{i,j\}\in E(G)
 \quad\Longleftrightarrow\quad
 \Delta_{ij}\text{ is a union of strata of }\Sg_d(G).
\]
Consequently, for simple graphs $G,H$ on $[n]$,
\[
 \Sg_d(G)=\Sg_d(H)\quad\Longleftrightarrow\quad E(G)=E(H).
\]
\end{theorem}

\begin{proof}
If $e=\{i,j\}$ is an edge, the stress positive only on $e$ belongs to $W_G(P)$
exactly when $p_i=p_j$: its only possibly nonzero forces are
$p_j-p_i$ and $p_i-p_j$. Thus a coordinate-sign-preserving homeomorphism of
fibres cannot identify a collapsed and a noncollapsed placement of $e$.
The collision diagonal is therefore a union of whole equivalence classes,
hence of strata. This elementary collision detector is already recorded in
Panina's Lemma~3.1 (preprint Section~3, Lemma~4) \cite{Panina}; we claim no novelty for it.

Conversely, suppose $\{i,j\}$ is not an edge of $G$. Choose $a\in\R^d\setminus
\{0\}$ and, for $-1/2\leq t\leq1/2$, set
\[
 p_k(t)=x_k(t)a,\qquad x_i(t)=t,\quad x_j(t)=0,\quad
 x_k(t)=k+1\ \ (k\notin\{i,j\}).
\]
The only pair that can collide on this path is $\{i,j\}$, so no $G$-edge
collapses. Orient the edges arbitrarily and let $B_G$ be the vertex-edge
incidence matrix, with $+1$ at the initial endpoint and $-1$ at the final one.
For an oriented edge $(u,v)$ put
$D_t=\operatorname{diag}(x_v(t)-x_u(t))$. Since $a\ne0$,
\[
 W_G(P(t))=\ker(B_GD_t),\qquad
 T_t=D_t^{-1}D_0,
 \qquad B_GD_tT_t=B_GD_0.
\]
Every entry of $D_t$ stays nonzero with constant sign. Hence $T_t$ is a
positive diagonal linear homeomorphism from $W_G(P(0))$ onto $W_G(P(t))$;
its inverse is $D_0^{-1}D_t$. The construction is independent of edge
orientation and preserves zero coordinates as well as both nonzero signs.
For an edgeless graph it is the unique map between zero stress spaces.

The entire continuous path thus lies in one fibre-equivalence class. Its
connected image lies in one connected component of that class, hence in
one stratum. But $P(0)\in\Delta_{ij}$ and $P(1/4)\notin\Delta_{ij}$.
The diagonal cannot be a union of whole strata. This proves the equivalence.
Identical partitions now recover identical edge sets; the reverse implication
is immediate. For $n=1$ there is only the empty simple graph.
\end{proof}

The collinear path is a witness \emph{inside} the full configuration space;
the theorem is not a classification restricted to collinear placements.
Stress-space dimension alone does not supply the argument: it is the labelled
coordinate support that detects a collapsed edge.

\begin{corollary}
A specified permutation of vertex labels carries $\Sg_d(G)$ onto
$\Sg_d(H)$ exactly when it carries $G$ onto $H$.
\end{corollary}
\begin{proof}
Reindexing points and edge weights carries the equilibrium equations for $G$
to those for the relabelled graph. Apply Theorem~\ref{thm:main}.
\end{proof}

\section{The original question and the comparison convention}
Karpenkov's Problem~4 \cite[p.~22]{Karpenkov} asks which subgraphs of $K_n$
define the same stratifications. Its preceding definitions use connected
strata in the full labelled Cartesian configuration space. Under literal
equality of these defined partitions, Theorem~\ref{thm:main} gives the complete
answer: each equality class contains exactly one labelled graph. The source
does not separately define a cross-graph abstract or ambient equivalence.
We do not infer an author's unwritten intention or claim to classify such
different comparisons, generic-only spaces, or collision-free spaces.

The distinction has a concrete mathematical effect. Let $G$ be the
four-vertex path with edges $12,23,34$, and $H$ the star with edges $12,13,14$.
For any forest, leaf elimination shows that stresses are supported exactly on
its collapsed edges; those coordinates are arbitrary and independent.
The invertible linear transformation
\[
 q_1=p_1,\quad q_2=p_2,\quad
 q_3=p_1+p_3-p_2,\quad q_4=p_1+p_4-p_3
\]
turns the path's three edge differences into the star's three edge differences.
It therefore maps their full stress-equivalence classes, and their connected
components, bijectively under that edge correspondence. The graphs have
different degree multisets, so they are nonisomorphic, although their
stratifications are ambient-linearly equivalent. Their partitions are unequal
in the fixed labelled space by the theorem.

\paragraph{Prior credit and verification.}
The stress-space convention comes from \cite{DKS,Karpenkov}; the
singleton-support detector is credited above. Collinear order invariance of
the stress matroid is also noted in \cite[Appendix~A.3; preprint v4, Section~4.1]{Panina}.
We give the incidence transport explicitly for a self-contained proof and
claim no novelty for these ingredients. The conclusion compared here is
reconstruction from the full fixed-label partition. Panina's affine-quotient and
noncollapsed-edge framework is not substituted for the full space used here.
The supplement records source and priority checks, independent adversarial
derivations, and reproducible exact-rational calculations. Finite tests
corroborate the identities and boundary cases; they do not prove the universal
or connectedness statements and provide no formal proof-assistant certificate.
No first-priority claim is made for the cited ingredients.

\paragraph{AI use and review status.}
AI tools were used extensively in developing and checking the arguments,
auditing sources, and preparing this manuscript and its supporting materials.
This is an unrefereed research note. The independent adversarial reviews
described in the supplement are AI reviews; conventional human peer review
or journal acceptance is not claimed.

\begin{thebibliography}{9}
\bibitem{DKS} F.~Doray, O.~Karpenkov and J.~Schepers,
\emph{Geometry of configuration spaces of tensegrities},
Discrete Comput. Geom. \textbf{43} (2010), 436--466.
\href{https://doi.org/10.1007/s00454-009-9229-4}{doi:10.1007/s00454-009-9229-4}.
Author preprint: \href{https://arxiv.org/abs/0806.4976v2}{arXiv:0806.4976v2}.
\bibitem{Karpenkov} O.~Karpenkov,
\emph{Open Problems on Configuration Spaces of Tensegrities},
Arnold Math. J. \textbf{4} (2018), 19--25.
\href{https://doi.org/10.1007/s40598-018-0080-7}{doi:10.1007/s40598-018-0080-7}.
\bibitem{Panina} G.~Panina,
\emph{A universality theorem for stressable graphs in the plane},
Ars Math. Contemp. \textbf{18} (2020), 137--148.
\href{https://doi.org/10.26493/1855-3974.641.e06}{doi:10.26493/1855-3974.641.e06}.
Published Lemma~3.1 and Appendix~A.3; author version: \href{https://arxiv.org/abs/1902.07212v4}{arXiv:1902.07212v4},
Section~3, Lemma~4.
\end{thebibliography}
\end{document}
